\subsection{2.4 - Análisis de Riesgos}

La gestión de riesgos constituye un proceso fundamental para asegurar el éxito del presente proyecto, permitiendo la identificación, análisis y mitigación proactiva de los factores que podrían comprometer sus objetivos. Dada la naturaleza dinámica de nuestro ``Sistema de Administración de Consorcios'', el enfoque adoptado se basa en un marco de trabajo continuo que abarca cuatro etapas críticas: la identificación de amenazas potenciales para el negocio, el producto y el cronograma del proyecto; el análisis cualitativo de su probabilidad e impacto; la planificación de estrategias de prevención y contingencia; y, finalmente, la supervisión constante para ajustar las medidas de control según la evolución del desarrollo.  

Este análisis permite al equipo anticipar posibles dificultades, optimizar la asignación de recursos y garantizar la calidad y fiabilidad del software final frente a incertidumbres tanto técnicas como organizacionales.

\begin{figure}[H]
    \centering
    \begin{tabularx}{\textwidth}{|p{5cm}|p{3cm}|X|}
    \hline
            \textbf{Riesgo} 
        
        & 
        
            \textbf{Tipo} 
        
        & 
        
            \textbf{Descripción} 
        
    \\ \hline
    
            \textbf{Error en cálculos financieros} 
        
        & 
        
            Producto 
        
        & 
        
            Errores en la lógica del prorrateo automático de expensas o en el cálculo de intereses por mora. 
        
    \\ \hline
    
            \textbf{Falla en seguridad de datos} 
        
        & 
        
            Producto 
        
        & 
        
            Acceso no autorizado o filtración de datos sensibles de los vecinos y registros del Libro de Visitas. 
        
    \\ \hline
    
                \textbf{Falla de integración de pagos} 
        
        & 
        
            Proyecto (producto) 
        
        & 
        
            Problemas técnicos al conectar el sistema con la pasarela de pagos para la generación de recibos. 
        
    \\ \hline
    
            \textbf{Nuevos requerimientos del cliente} 
        
        & 
        
            Negocio 
        
        & 
        
            Solicitud de cambios o nuevas funciones (ej. agregar un módulo de votación a mitad del desarrollo). 
        
    \\ \hline
    
            \textbf{Baja adopción tecnológica} 
        
        & 
        
            Negocio 
        
        & 
        
            Resistencia de los vecinos (especialmente adultos mayores) a utilizar la aplicación para reservas o pagos. 
        
    \\ \hline
    
            \textbf{Falta de experiencia técnica} 
        
        & 
        
            Proyecto 
        
        & 
        
            Dificultad para manejar las herramientas de desarrollo elegidas para el módulo de notificaciones push. 
        
    \\ \hline
    
            \textbf{Interfaz poco intuitiva} 
        
        & 
        
            Producto 
        
        & 
        
            Diseño de pantallas confuso que dificulte la carga de tickets de mantenimiento por parte del vecino. 
        
    \\ \hline
    
            \textbf{Ausencia de integrantes} 
        
        & 
        
            Proyecto 
        
        & 
        
            Inasistencia de un miembro del equipo por razones de salud o personales durante las fechas de entrega. 
        
    \\ \hline
    \end{tabularx}
    \caption{Tabla de los riesgos que se pensaron}
\end{figure}

Tras la identificación, se procede a la valoración cualitativa de cada riesgo (Probabilidad y Seriedad). Finalmente, se define el plan de gestión y la estrategia de mitigación para cada uno.

\begin{figure}[H]
    \centering
    \begin{tabularx}{\textwidth}{|X|l|l|}
    \hline
            \textbf{Riesgo} 
        
        & 
        
            \textbf{Probabilidad} 
        
        & 
        
            \textbf{Seriedad} 
        
    \\ \hline
    
            \textbf{Falla de seguridad de datos} 
        
        & 
        
            Baja 
        
        & 
        
            Alta 
        
    \\ \hline
    
            \textbf{Error en cálculos financieros} 
        
        & 
        
            Media 
        
        & 
        
            Alta 
        
    \\ \hline
    
            \textbf{Ausencia de integrantes} 
        
        & 
        
            Media 
        
        & 
        
            Alta 
        
    \\ \hline
    
            \textbf{Nuevos requerimientos del cliente} 
        
        & 
        
            Alta 
        
        & 
        
            Media 
        
    \\ \hline
    
            \textbf{Baja adopción tecnológica} 
        
        & 
        
            Media 
        
        & 
        
            Media 
        
    \\ \hline
    
            \textbf{Falla de integración de pagos} 
        
        & 
        
            Alta 
        
        & 
        
            Media 
        
    \\ \hline
    
            \textbf{Falta de experiencia técnica} 
        
        & 
        
            Media 
        
        & 
        
            Baja 
        
    \\ \hline
    
            \textbf{Interfaz poco intuitiva} 
        
        & 
        
            Baja 
        
        & 
        
            Baja 
        
    \\ \hline
    \end{tabularx}
    \caption{Tabla con las probabilidades y seriedad de los riesgos}
\end{figure}

A continuación, se detalla la estrategia de gestión y el plan de acción específico para la prevención o mitigación de cada riesgo.

\begin{figure}[H]
    \centering
    \begin{tabularx}{\textwidth}{|p{5.5cm}|p{2.5cm}|X|}
    \hline
        \textbf{Riesgo} 
        
        & 
        
            \textbf{Estrategia} 
        
        & 
        
            \textbf{Plan de Gestión (Detalle)} 
        
    \\ \hline
    
            \textbf{Falla de seguridad de datos} 
        
        & 
        
            Prevención 
        
        & 
        
            Implementar cifrado de datos sensibles y auditorías de acceso desde la fase de diseño. 
        
    \\ \hline
    
            \textbf{Error en cálculos financieros} 
        
        & 
        
            Prevención 
        
        & 
        
            Realizar pruebas unitarias intensivas sobre la lógica del prorrateo y validaciones cruzadas. 
        
    \\ \hline
    
            \textbf{Ausencia de integrantes} 
        
        & 
        
            Minimización 
        
        & 
        
            Documentar todo el código y procesos para que cualquier miembro pueda retomar tareas de otro. 
        
    \\ \hline
    
            \textbf{Nuevos requerimientos del cliente} 
        
        & 
        
            Prevención 
        
        & 
        
            Establecer un alcance claro (Scope) y utilizar un mecanismo formal de gestión de cambios. 
        
    \\ \hline
    
            \textbf{Baja adopción tecnológica} 
        
        & 
        
            Minimización 
        
        & 
        
            Crear tutoriales breves y realizar una prueba de usabilidad con un grupo control de adultos mayores. 
        
    \\ \hline
    
            \textbf{Falla de integración de pagos} 
        
        & 
        
            Contingencia 
        
        & 
        
            Tener un plan de respaldo manual para la generación de recibos si la API de pago falla. 
        
    \\ \hline
    
            \textbf{Falta de experiencia técnica} 
        
        & 
        
            Prevención 
        
        & 
        
            Capacitación temprana mediante workshops técnicos o tutoriales sobre notificaciones push. 
        
    \\ \hline
    
            \textbf{Interfaz poco intuitiva} 
        
        & 
        
            Prevención 
        
        & 
        
            Realizar prototipado rápido y validación con usuarios reales antes de la codificación final. 
        
    \\ \hline
    \end{tabularx}
    \caption{Estrategias para cada riesgo}
\end{figure}
